\documentclass[12pt]{article}
\usepackage[utf8]{inputenc}
\usepackage[T1]{fontenc}
\usepackage{hyperref}
\title{Tagoi Language Resources for K-12 Curriculum Development}
\author{Compiled by Claude Sonnet 5 (claude-sonnet-5) for Andr\'as Kornai}
\date{}
\begin{document}
\maketitle

\section{Introduction}
Tagoi (ISO 639-3: tag; Glottocode: tago1246; also Tagoy, Moreb, Tumale) is a Rashad-cluster language of the Nuba Mountains, South Kordofan, Sudan. A figure of 13,000 speakers (SIL 1982) is cited in search results. Its classification as Kordofanian has been questioned, but it is accepted as Niger-Congo. Very little exists beyond linguistic survey work.

\section{Primary Educational Resources}
No primers, textbooks or school materials were found in search results.

\section{Linguistic and Reference Materials}
\begin{itemize}
\item A preliminary linguistic survey of the Tagoi language, Nuba Mountains, Sudan (ELDP-funded, 2007): \url{https://cultureincrisis.org/projects/a-preliminary-linguistic-survey-of-the-tagoi-language-nuba-mountains-sudan}
\item \emph{Nuba Mountain Language Studies: Contemporary Perspectives} (K\"oln; includes Tagoi chapters per search summary): \url{https://monographs.ub.uni-koeln.de/usbk/catalog/book/38}
\item Glottolog entry: \url{https://glottolog.org/resource/languoid/iso/tag}
\item Endangered Languages Project entry: \url{https://endangeredlanguages.com/node/84960}
\end{itemize}

\section{Religious and Community Resources}
Joshua Project lists the Tagoy people group (population data only; no language materials seen): \url{https://joshuaproject.net/people_groups/15178}

\section{Digital and Audio Resources}
None found in search results.

\section{Curriculum Development Context}
No written standard, school programme or teaching materials were found. Any curriculum would start from documentation and community work, with attention to the security situation in Sudan. Arabic is the regional language of education.
\end{document}